\scanpage{79}
% archive foliation: 78
On est obligé de renforcer la propriété universelle de $\otimes$ et $\Gamma^p$, pour avoir les propriétés d'assoc. etc qu'on souhaite.

Soit $X\times Y\longrightarrow\deleted{\unsure{\unsure{\mathcal P|P}|Q}}T$ un morphisme de bidegré $(1,1)$, alors pour tt $Z$, on a
\[
X\times Y\times Z\longrightarrow T\times Z
\]
d'où pour tt $M$
\[
\unsure{\mathcal P|P}_{1,n}(T,Z;M)\longrightarrow \unsure{\mathcal P|P}_{1,1,n}(X,Y,Z;M).
\]
On dit que $T$ est un produit \uline{tensoriel strict} de $X$ et $Y$ si l'hom précédent est un iso \deleted{\unsure{et alors $X\otimes Y=T$}} pour tt $n$ et tout $M$. Faisons $n=0$, $Z=0$, on trouve
\[
\unsure{\mathcal P|P}_{1,\deleted{\unsure{0}}}(T,M)\simeq \unsure{\mathcal P|P}_{1,1}(X,Y;M)
\]
donc $T$ est un produit tensoriel. C'est en exigeant qu'on a des $\otimes$ \uline{stricts} \deleted{que} \inserted{\unsure{définis}} \inserted{\unsure{pour des facteurs} (\deleted{\unsure{qui}} \unsure{il faut que la condition précédente} \unsure{soit prise} \unsure{pour un nombre} \unsure{quelconque de facteurs})} qu'on obtient une structure \unsure{sur $\unsure{\mathcal P|P}_1$} ass. comm. compatible avec
\[
\deleted{\unsure{\operatorname{Hom}_{\unsure{\mathcal P|P}_1}}}\quad \unsure{\mathcal P|P}_1(X_1\otimes\cdots\otimes X_r;M)\simeq \unsure{\mathcal P|P}_{1,\ldots,1}(X_1,\ldots,X_r;M).
\]
On sait qu'un objet unité (\unsure{sous réserve d'une|sans recours à une} donnée d'associativité) si elle existe, est unique. \margin{?} Il faut examiner pourquoi ça devrait être $\Gamma^0(0)$.

Revenons aux $\Gamma^n$, soit $X=X_1\times\cdots\times X_r$, tel que $\Gamma^n(X)$ existe. Considérons les sous foncteurs $\unsure{\unsure{\mathcal P|P}|\Gamma}_{\nu=(\nu_1,\ldots,\nu_r)}(X_1,\ldots,X_r;M)$ du foncteur \inserted{$\unsure{\mathcal P|P}_1\longrightarrow(\operatorname{Ab})$} (can. représenté par $\Gamma^n(X)$), supposons qu'ils soient aussi \unsure{cor} (automatique si $\unsure{\mathcal P|P}_1$ est karoubienne) alors $\Gamma^{\nu_1,\ldots,\nu_r}(X_1,\ldots,X_r)=\Gamma^\nu(X_1,\ldots,X_r)$ l'\deleted{\unsure{produit}} \inserted{\unsure{objet}} qui le représente. \deleted{\unsure{qui}}
(On va définir \unsure{sans réserve} d'existence des $\Gamma^{\nu_i}X_i$ \deleted{\unsure{produits}} \inserted{et des} tensoriels, \unsure{ou} dans $\check \unsure{\mathcal P|P}_1=\widehat{\unsure{\mathcal P|P}_1^\circ}{}^\circ$ si

